\section{Conclusion}
\label{sec:conclusion}

This paper presents \benchmark, the final generated dataset of Java classes and method-reachability questions, together with the controlled algorithm used to construct it. Each question has a Boolean label and an executable-graph oracle, allowing decisions and supporting traces to be evaluated independently.

The evaluation shows that correct decisions often lack valid executable evidence. Conditions, passed values, and the exhaustive exploration required by negative questions produce the strongest difficulties. Trace diagnostics locate errors around disabled calls, omitted methods, and longer value dependencies.

Providing the \texttt{main} entry point and explicitly describing the deterministic execution improves the observed scores, but the difficulties with guarded value propagation remain. The main comparison also increases the reasoning budget. At a fixed budget, Qwen3.8-Flash-Next improves its trace scores on the difficult subset while producing more non-responses. These findings show why reachability should be evaluated through both decisions and executable evidence, with the prompt context made explicit. Further evaluation on more models and real programs would test how broadly these difficulties transfer.
